%!TEX root = ../diploma.tex

%Отступ в первом абзаце раздела
\usepackage{indentfirst}
%Точка после номера раздела
\usepackage{misccorr}
% Настройка списков с пользовательскими метками (label=...)
\usepackage{enumitem}

\usepackage{hyperref}
%Тонкое оформление перекрёстных ссылок
\usepackage[russian]{cleveref}
\usepackage[dvipsnames, table, hyperref]{xcolor}
\usepackage{csquotes}
\usepackage{hyphenat}
% \hyphenation{ма-те-ма-ти-ка вос-ста-нав-ли-вать}


% Принудительное уменьшение размера заголовков глав
\usepackage{titlesec}

% Уменьшение размера заголовков секций
\titleformat{\chapter}[hang]{\large\bfseries}{\thechapter}{1em}{}
\titleformat{\section}{\large\bfseries}{}{0pt}{}
\titleformat{\subsection}{\normalsize\bfseries}{}{0pt}{}
\titleformat{\subsubsection}{\normalsize\itshape}{}{0pt}{}
% Для ненумерованных глав (звёздочка)
\titleformat{name=\chapter,numberless}[hang]{\large\bfseries}{}{0pt}{} % Если вы хотите ещё меньший шрифт, замените \large на \normalsize или \small.
% Настройка отступов до и после заголовка
\titlespacing*{\chapter}{0pt}{0pt}{20pt}

% Скрываем номера section/subsection/subsubsection в оглавлении
\makeatletter
\let\origl@section\l@section
\renewcommand*{\l@section}[2]{%
  \begingroup
  \let\numberline\@gobble
  \origl@section{#1}{#2}%
  \endgroup
}
\let\origl@subsection\l@subsection
\renewcommand*{\l@subsection}[2]{%
  \begingroup
  \let\numberline\@gobble
  \origl@subsection{#1}{#2}%
  \endgroup
}
\let\origl@subsubsection\l@subsubsection
\renewcommand*{\l@subsubsection}[2]{%
  \begingroup
  \let\numberline\@gobble
  \origl@subsubsection{#1}{#2}%
  \endgroup
}
\makeatother
